The overlap table gives the original parent row number, not its index in the reordered construction array.
\begin{center}
\begin{tabular}{@{}rrrr@{}}
\toprule Target index $j$ & $H_j$ & Original row & Sign\\\midrule
0 & 358 & 14 & $+1$\\
1 & 608 & 7 & $+1$\\
2 & 939 & 21 & $+1$\\
3 & 1147 & 10 & $+1$\\
4 & 1432 & 8 & $+1$\\
5 & 1739 & 17 & $+1$\\
6 & 1989 & 16 & $+1$\\
7 & 2170 & 6 & $+1$\\
8 & 2512 & 1 & $+1$\\
9 & 2699 & 13 & $+1$\\
10 & 2972 & 4 & $+1$\\
11 & 3305 & 2 & $+1$\\
12 & 3492 & 19 & $+1$\\
13 & 3683 & 11 & $+1$\\
14 & 3914 & 20 & $+1$\\
15 & 4108 & 5 & $+1$\\
16 & 4399 & 18 & $+1$\\
17 & 4697 & 15 & $+1$\\
18 & 4926 & 22 & $+1$\\
19 & 5224 & 3 & $+1$\\
20 & 5469 & 12 & $+1$\\
21 & 5775 & 9 & $-1$\\
\bottomrule
\end{tabular}
\end{center}
The uniform placement uses all positive signs and original row order $1,2,3,4,6,7,\ldots,22,5$, with $H_j=358+359j$. Both layouts use the weights in equation~\eqref{eq:weights}.
